% ---------------- ABSTRAK (Indonesia) ----------------
\clearpage
\phantomsection
\addcontentsline{toc}{chapter}{ABSTRAK}
\begin{center}
    \textbf{ABSTRAK}\\[0.5cm]
\end{center}

Pembangkit Listrik Tenaga Surya (PLTS) terapung beroperasi pada lingkungan yang
dipengaruhi air dan kelembapan sehingga modul fotovoltaik (PV) berpotensi
mengalami kondisi abnormal, termasuk kondisi terendam. Modul yang pernah
terendam tidak dapat langsung dinyatakan layak hanya berdasarkan inspeksi
visual sehingga diperlukan evaluasi berbasis parameter elektrik. Laporan kerja
praktik ini membahas analisis kelayakan modul PV pasca-terendam pada PLTS
Terapung Cirata menggunakan pengujian kurva arus--tegangan (I--V) dan indikator
arus--tegangan. Analisis dilakukan terhadap 59 data pengujian modul tipe
JKM560N-72HL4-BDV yang diuji secara \textit{onshore} menggunakan \textit{I--V
curve tracer} dengan koreksi parameter ke \textit{Standard Test Conditions}
(STC). Evaluasi mencakup perbandingan parameter STC terhadap \textit{nameplate},
perhitungan Fill Factor, serta analisis indikator $R_{VM}$ dan $R_{IM}$ untuk
menyusun interpretasi awal terhadap kegiatan \textit{operation and maintenance}.
Hasil menunjukkan 6 data (10,17\%) berstatus Ok dan 53 data (89,83\%) berstatus
Not Ok, dengan rata-rata daya maksimum STC kelompok Not Ok sekitar 501,6 W atau
10,4\% di bawah nilai \textit{nameplate} 560 W. Indikator $R_{VM}$ memisahkan
kedua kelompok lebih jelas dibandingkan $R_{IM}$ sehingga perbedaan performa
cenderung dominan pada sisi tegangan. Karena sebagian data Ok diperoleh pada
irradiansi rendah, status kelayakan perlu ditafsirkan secara hati-hati dan
dikonfirmasi melalui pengujian ulang pada kondisi irradiansi yang lebih
representatif serta melalui pengujian tahanan isolasi (\textit{insulation
resistance test}) sebelum keputusan pemasangan kembali dilakukan.

\vspace{1em}
\noindent Kata kunci: PLTS terapung, modul PV terendam, pengujian I--V,
indikator arus--tegangan, Fill Factor